% ==============================================================
\section{実験}\label{sec:experiments}
% ==============================================================

\subsection{実験設計}

\subsubsection{合成データ生成}

ベース語彙200アイテム（各アイテムの出現確率 $p_{\text{base}} = 0.03$），
$N = 5{,}000$ トランザクションの合成データを生成する．
外部イベントの効果は\emph{植え込み信号}としてモデル化する．
3種類の信号を1つのデータセットに混在させることで，
パイプラインの弁別能力を多面的に評価する:

\noindent\textbf{Type A}（語彙内ブースト）:
ベース語彙内のアイテムペアを使用し，
イベント期間外にも $p_{\text{base}}$ でベースラインサポートを持たせる．
イベント活性期間中に確率 $\beta$ で全パターンアイテムを同時挿入し
サポートを上昇させる．
magnitude単独では不十分な\textbf{現実的条件}を構成する．

\noindent\textbf{Type B}（イベント非対応密集パターン）:
特定期間にサポートが上昇するが，いずれのイベントとも
対応しないパターン．季節性やトレンドを模擬し，
パイプラインが\textbf{無関係な変動を棄却}できるかを検証する．

\noindent\textbf{Type C}（デコイイベント）:
パターン変動と無関係なイベント．

\subsubsection{評価指標}

\begin{itemize}
\item \textbf{P/R/F1}: 予測と正解の厳密一致に基づく
    Precision，Recall，F1
\item \textbf{偽帰属率 (FAR)}: Type Bパターンがイベントに誤帰属された割合
\end{itemize}

\subsubsection{デフォルト設定}

表~\ref{tab:default_params}にデフォルトパラメータを示す
（各パラメータの選択根拠は付録~\ref{app:params}参照）．

\begin{table}[t]
\caption{デフォルトパラメータ}
\label{tab:default_params}
\centering
\begin{tabular}{ll}
\toprule
パラメータ & 値 \\
\midrule
ウィンドウサイズ $W$ & 50 \\
最小サポート $\theta$ & 3 \\
帰属対象パターン & $|P| \geq 2$\footnotemark[1] \\
最大パターン長 $L$ & 制約なし\footnotemark[2] \\
変化点検出法 & 閾値交差 \\
$\sigma$ & $W$ (= 50) \\
置換回数 $B$ & 5{,}000 \\
有意水準 $\alpha$ & 0.10 \\
補正方法 & BH \\
変化量正規化 & 平方根（式~(\ref{eq:mag_sqrt})） \\
最小サポート振幅 $\delta$ & 5 \\
重複排除 & 有効 \\
\bottomrule
\end{tabular}
\end{table}

\footnotetext[1]{%
本手法は共起パターンのサポート変動を分析対象としており，
単一アイテム（$|P|=1$）の頻度変動は通常の時系列分析で捕捉可能なため
帰属対象外とする（定義~\ref{def:attribution}参照）．}
\footnotetext[2]{%
Aprioriの反単調性により，サポート閾値を満たすパターンが存在しなくなった時点で
探索が自然に停止するため，最大パターン長に明示的な制約は設けない．
本実験データでは実質的に$L \leq 4$で停止する．}

\subsubsection{評価方針}

本研究が定式化するイベント帰属問題（定義~\ref{def:attribution}）は，
数千のパターンサポート時系列を複数の外部イベントに同時に帰属させ，
多重検定を制御しつつ副次パターンの偽陽性を排除するタスクである．
表~\ref{tab:related_comparison}に示した通り，
既存手法（ITS，CausalImpact，ECA等）はいずれもこの問題を直接解くものではなく，
入出力の構造が本質的に異なるため，同一条件での定量比較は公正に行えない．
したがって本実験では，外部手法との比較ではなく，
パイプラインの各構成要素（置換検定，BH補正，重複排除，変化量正規化）を
段階的に除去する\textbf{アブレーション分析}（EX2）により，
各ステップの必要性と寄与を実証する．

% ------------------------------------------------------------------
\subsection{EX1: コア帰属精度}
\label{sec:exp_ex1}

\textbf{RQ}: パイプラインはイベント起因の変動を正しく帰属し，
イベント無関係の密集パターンを棄却できるか？
多様なデータ特性（信号強度，構造条件，アイテム頻度分布）に対する頑健性を評価する．

\paragraph{信号強度}
Type A（語彙内ブースト3パターン），
Type B（イベント非対応2パターン），
Type C（デコイ2件）を混在させ，
$\beta \in \{0.1, 0.2, 0.3, 0.5\}$ で5 seed実行した
（表~\ref{tab:ex1}）．

\begin{table}[t]
\caption{EX1a 結果: 信号強度別の帰属精度（5 seeds平均）}
\label{tab:ex1}
\centering
\begin{tabular}{ccccc}
\toprule
$\beta$ & Precision & Recall & F1 & FAR \\
\midrule
0.1 & 0.74 & 1.00 & 0.84 & 0.20 \\
0.2 & 0.74 & 1.00 & 0.84 & 0.10 \\
0.3 & 0.62 & 0.80 & 0.69 & 0.10 \\
0.5 & 0.45 & 0.60 & 0.51 & 0.00 \\
\bottomrule
\end{tabular}
\end{table}

$\beta \leq 0.2$ の弱い信号でRecall = 1.00, F1 $\geq$ 0.84を達成している．
$\beta = 0.5$ ではF1 = 0.51とやや低下するが，
これは強いブーストにより真のパターンの上位集合も高サポートとなり，
重複排除で残存する副次パターンが増加するためである．
FARは0.00--0.20の範囲で推移し，
信号が強いほどType Aパターンの帰属スコアが高くなり，
相対的にType Bの誤帰属が抑制される．

\paragraph{構造条件}
デフォルト設定（$\beta = 0.3$）で4つの構造条件を評価した結果を
表~\ref{tab:ex1_struct}に示す．

\begin{table}[t]
\caption{EX1b 結果: 構造条件別精度（5 seeds平均）}
\label{tab:ex1_struct}
\centering
\begin{tabular}{lcccc}
\toprule
条件 & Precision & Recall & F1 & FAR \\
\midrule
OVERLAP  & 0.59 & 0.80 & 0.67 & 0.00 \\
CONFOUND & 0.65 & 0.73 & 0.68 & 0.10 \\
DENSE    & 0.57 & 0.90 & 0.69 & 0.20 \\
SHORT    & 0.78 & 0.93 & 0.85 & 0.30 \\
\bottomrule
\end{tabular}
\end{table}

OVERLAPではF1 = 0.67, FAR = 0.00を達成し，
イベント区間が重複する条件でも誤帰属を完全に棄却できている．
CONFOUNDではType Bパターンがイベント近傍に配置される交絡条件にもかかわらず
FAR = 0.10と低く，置換検定が交絡の影響を適切に排除している．
DENSEではFAR = 0.20と上昇するが，Recall = 0.90は維持されている．
SHORTでは短期間イベント（持続期間80）でF1 = 0.85と最も高い精度を達成した．

\paragraph{アイテム頻度分布}
EX1a--EX1bでは全アイテムが一様な出現確率 $p_{\text{base}} = 0.03$ を持つ．
しかし実際のトランザクションデータではアイテム頻度は
Zipf（べき乗則）分布に従うことが知られている~\cite{fournier2017survey}．
高頻度アイテムは自然に多くのパターンに含まれるため，
帰属スコアや重複排除の有効性に影響しうる．

アイテム $k$（$k = 1, \ldots, 200$）の出現確率を
$p(k) = \min(C / k^{\alpha_z},\; 0.10)$ とし（$C$ は中央値アイテムの確率が
0.03 程度になるよう正規化，最大確率を0.10で打ち切り），
EX1a と同一の信号構造を使用する．
最小サポートは $\theta = 5$ とし，Zipf 由来の高頻度パターンを適度にフィルタする．
3条件を評価する:
Zipf-1.0（$\alpha_z = 1.0$，標準 Zipf），
Zipf-1.5（$\alpha_z = 1.5$，重裾分布），
Correlated（Zipf-1.0 + 相関アイテムペア，ヘッドアイテム間の共起確率 0.5--0.7）．
結果を表~\ref{tab:ex1_zipf}に示す．

\begin{table}[t]
\caption{EX1c 結果: アイテム頻度分布別の精度（5 seeds平均）}
\label{tab:ex1_zipf}
\centering
\begin{tabular}{lcccc}
\toprule
条件 & Precision & Recall & F1 & FAR \\
\midrule
一様分布（$\beta=0.3$） & 0.62 & 0.80 & 0.69 & 0.10 \\
Zipf-1.0 & 0.32 & 1.00 & 0.48 & 0.10 \\
Zipf-1.5 & 0.24 & 1.00 & 0.38 & 0.10 \\
Correlated & 0.17 & 0.53 & 0.25 & 0.10 \\
\bottomrule
\end{tabular}
\end{table}

一様分布からZipf分布に移行すると精度が低下する．
Zipf 分布下では高頻度アイテムペアの自然な共起変動がノイズ仮説を生成し，
Precision を低下させることが原因である．
Correlated 条件では相関アイテムペアが定常的に高いサポートを持つため，
イベント起因の変動がノイズに埋没しやすく，
Recall も低下する．
平方根正規化の効果はEX2のアブレーション分析で検証する．

\paragraph{帰無条件検証}
植え込み信号を一切含まない帰無データ（EX1と同一の生成条件，
5件のランダムイベント，持続期間300）で20 seed実行した結果，
全seedにおいて有意帰属数は0であった．
BH手続きと円形シフト置換検定の組み合わせが
帰無仮説下のp値を十分に大きく保ち，
過剰な偽発見を生じないことが経験的に確認された．

% ------------------------------------------------------------------
\subsection{EX2: アブレーション分析}
\label{sec:exp_ex2}

\textbf{RQ}: スコア構成要素，変化量正規化，およびパイプラインの各ステップは精度にどう寄与するか？

\paragraph{スコア構成要素}

各構成要素の必要性を実証する2つのシナリオを設計し，
全3アブレーションモードを各シナリオに適用する（2シナリオ $\times$ 3モード $\times$ 5 seeds = 30条件）:

\noindent\textbf{Scenario A} (prox必要):
パターン $(5, 15)$ に2つの密集区間 --- 1つはイベント近傍（因果的），
1つは遠方（偶発的）．proxなしでは両者に等しいスコアが付与される．

\noindent\textbf{Scenario B} (mag必要):
パターン $(5, 15)$ が2イベントの影響を受けるが，
ブースト量が $\beta=0.5$ と $\beta=0.1$ で大小異なる．
magなしでは両者が同等に扱われる．

結果を表~\ref{tab:ex2}に示す．

\begin{table}[t]
\caption{EX2a 結果: シナリオ別アブレーション（5 seeds平均F1）}
\label{tab:ex2}
\centering
\begin{tabular}{lcc}
\toprule
構成 & Scenario A & Scenario B \\
\midrule
Full: $\text{prox} \times \text{mag}$ & 0.60 & 0.53 \\
mag のみ（prox除去） & 0.00 & 0.00 \\
prox のみ（mag除去） & 0.60 & 0.28 \\
\bottomrule
\end{tabular}
\end{table}

表~\ref{tab:ex2}から，\textbf{proximity が帰属に不可欠}であることが明らかである．
本パイプラインではハード距離カットオフ（max\_distance）を使用せず，
prox の指数減衰（式~(\ref{eq:proximity})）が距離に応じた自然な重み付けを行う．
このため，mag のみ（prox除去）の構成は時間的対応付け情報を一切持たず，
全変化点--イベントペアに等しいスコアが付与される結果，
\textbf{両シナリオでF1 = 0.00}に崩壊する．
prox のみ（mag除去）ではScenario Aで F1 = 0.60，Scenario Bで F1 = 0.28 と部分的に機能するが，
変化量情報の欠如により精度は不十分である．

\paragraph{構成要素間の相補性}
Full構成（$\text{prox} \times \text{mag}$）が
Scenario AでF1 = 0.60，Scenario BでF1 = 0.53を達成する一方，
いずれの単独構成要素もすべてのシナリオで同等の精度を達成できない．
特に，prox除去時の完全な精度崩壊（F1 = 0.00）は，
時間的近接度が帰属の\textbf{必要条件}であることを示す．
これは乗算構成がいずれかの構成要素が0に近ければスコア全体を抑制する
排他的フィルタとして機能し，偽陽性の抑制に寄与するためである．
平方根正規化下ではFull構成のF1がやや低下しているが，
これはベースラインサポートによる除算が弁別力を一部相殺するためであり，
Zipf的な分布を含む多様な条件への頑健性とのトレードオフと理解できる．

\paragraph{パイプラインステップの寄与}
ナイーブベースライン（Steps 1--3のみ: 変化点検出 + スコアリング，
閾値以上の候補をすべて出力）と，パイプラインの各ステップを
段階的に追加した3つの変種を比較する:

\begin{itemize}
\item \textbf{Naive}: Steps 1--3のみ（検定なし）
\item \textbf{+PermTest}: per-pattern 置換検定を追加（グローバル補正なし）
\item \textbf{+BH}: グローバルBH補正を追加
\item \textbf{Full}: Union-Find重複排除を追加（= 完全パイプライン）
\end{itemize}

EX1のデータ（$\beta = 0.3$, CONFOUND, DENSE）を使用し，
各条件 $\times$ 5 seedsで評価した（表~\ref{tab:ex2b}）．

\begin{table}[t]
\caption{EX2b 結果: ナイーブベースラインとの段階的比較（5 seeds平均）}
\label{tab:ex2b}
\centering
\begin{tabular}{llrrrrr}
\toprule
条件 & 手法 & P & R & F1 & FAR & \#Pred \\
\midrule
\multirow{4}{*}{$\beta$=0.3}
  & Naive     & 0.01 & 1.00 & 0.01 & 0.80 & 527 \\
  & +PermTest & 0.01 & 1.00 & 0.01 & 0.30 & 465 \\
  & +BH       & 0.29 & 0.80 & 0.39 & 0.20 & 9 \\
  & Full      & \textbf{0.62} & 0.80 & \textbf{0.69} & \textbf{0.10} & 3 \\
\midrule
\multirow{4}{*}{CONFOUND}
  & Naive     & 0.00 & 1.00 & 0.01 & 1.00 & 673 \\
  & +PermTest & 0.00 & 1.00 & 0.01 & 1.00 & 609 \\
  & +BH       & 0.13 & 1.00 & 0.22 & 0.90 & 25 \\
  & Full      & \textbf{0.65} & 0.73 & \textbf{0.68} & \textbf{0.10} & 4 \\
\midrule
\multirow{4}{*}{DENSE}
  & Naive     & 0.01 & 1.00 & 0.01 & 0.95 & 969 \\
  & +PermTest & 0.01 & 1.00 & 0.01 & 0.65 & 925 \\
  & +BH       & 0.30 & 0.97 & 0.39 & 0.40 & 31 \\
  & Full      & \textbf{0.57} & 0.90 & \textbf{0.69} & \textbf{0.20} & 10 \\
\bottomrule
\end{tabular}
\end{table}

Naiveは全条件でRecall = 1.00を達成するが，
Precision が極めて低く偽陽性を大量に出力する（F1 $\leq$ 0.01）．
スコア閾値のみでは偽陽性の制御が不可能であることを示す．

\textbf{+PermTest}（per-pattern置換検定）は出力数をわずかに削減するが，
パターン単位の補正ではグローバルな多重検定の負担に対応できず，
FPの削減効果は限定的である．

\textbf{+BH}（グローバルBH補正）が最大の偽陽性削減効果を持つ．
BH手続きがグローバルなp値順位に基づきFDRを制御することで，
個別パターンの検定では見逃される弱い偽陽性を効果的に棄却する．

\textbf{Full}（+重複排除）は残存する冗長パターン
（同一イベントに帰属されたアイテム重複・部分集合パターン）を除去し，
交絡パターンの誤帰属を効果的に抑制する．

\paragraph{変化量正規化の効果}
デフォルトの平方根正規化（式~(\ref{eq:mag_sqrt})）を除去した場合の
精度変化を，一様分布（$\beta = 0.3$）とZipf分布（Zipf-1.0）の
2条件で評価する（表~\ref{tab:ex2_norm}）．

\begin{table}[t]
\caption{EX2c 結果: 変化量正規化のアブレーション（5 seeds平均）}
\label{tab:ex2_norm}
\centering
\begin{tabular}{llcccc}
\toprule
条件 & 正規化 & P & R & F1 & FAR \\
\midrule
\multirow{2}{*}{一様（$\beta=0.3$）}
  & $\sqrt{}$正規化（デフォルト） & 0.62 & 0.80 & 0.69 & 0.10 \\
  & 正規化なし & 0.87 & 0.87 & 0.82 & 0.10 \\
\midrule
\multirow{2}{*}{Zipf-1.0}
  & $\sqrt{}$正規化（デフォルト） & 0.25 & 0.87 & 0.39 & 0.10 \\
  & 正規化なし & 0.19 & 0.60 & 0.29 & 0.10 \\
\bottomrule
\end{tabular}
\end{table}

一様分布では正規化なしの方が高い精度を示す（F1: 0.82 vs 0.69）．
これは全パターンのベースラインサポートが同程度であるため，
平方根正規化がスコアの弁別力をわずかに低下させることによる．
一方，Zipf分布では平方根正規化によりF1が0.29から0.39へ34\%改善する．
高頻度アイテムペアの絶対変化量が自然変動で大きくなりやすいところ，
正規化によりスコアが抑制され，真の信号が相対的に浮上するためである．
実世界のトランザクションデータは一般にZipf的な頻度分布を持つため，
平方根正規化をデフォルトとすることが妥当である．
アイテム頻度分布が既知で一様に近い場合は，
正規化なしに切り替えることで精度が向上しうる．

% ------------------------------------------------------------------
\subsection{EX3: 実データ検証}
\label{sec:exp_ex3}

\textbf{RQ}: 実世界データで意味のある結果を生成できるか？

Dunnhumby ``The Complete Journey'' データセット~\cite{dunnhumby2014complete}の
276{,}484件の購買バスケットと30件の販促キャンペーンを用いた探索的分析．
キャンペーンの開始・終了日は既知だが正解帰属は不明のため，
\textbf{定性的検証}として位置づける．

$W = 300$, $\theta = 5$, $B = 5{,}000$, $\alpha = 0.10$（BH補正），
$\delta = 3$, 重複排除有効．

\begin{table}[t]
\caption{EX3 結果: Dunnhumby 実キャンペーンデータ}
\label{tab:ex3_dunnhumby}
\centering
\begin{tabular}{lr}
\toprule
指標 & 値 \\
\midrule
トランザクション数 & 276{,}484 \\
パターン数（$|P| \geq 2$） & 1{,}006 \\
キャンペーン数 & 30 \\
有意帰属数（重複排除後） & 114 \\
帰属のあるキャンペーン数 & 30/30 \\
キャンペーンあたり平均帰属数 & 3.8 \\
合計処理時間 & 1{,}292 s \\
\bottomrule
\end{tabular}
\end{table}

\textbf{キャンペーンカバレッジ}: 全30キャンペーンに少なくとも1件の
有意な帰属が生成された（カバレッジ100\%）．
キャンペーンあたりの帰属数は平均3.8件であり，
分析者にとって管理可能な量である．

\paragraph{クーポン整合性による部分的検証}
Dunnhumby データセットには各キャンペーンのクーポン対象商品リスト（coupon.csv）が
付属しており，帰属結果の部分的な正解検証が可能である．
帰属されたパターン $(a, b)$ の少なくとも1つのアイテムが
当該キャンペーンのクーポン対象商品に一致する場合を
\textbf{クーポン整合的}（coupon-consistent）と定義する（表~\ref{tab:ex3_coupon}）．
この定義は1-item overlapに基づく緩い基準であり，
パターン全体の一致（2-item完全一致）ではクーポン対象商品リストの
カバレッジが限定的すぎるため適用困難である．
したがって，クーポン整合率は帰属の妥当性の\textbf{上界的な部分指標}であり，
整合的であっても偶然一致の可能性を排除しない．

\begin{table}[t]
\caption{EX3: クーポン整合性分析（キャンペーンタイプ別）}
\label{tab:ex3_coupon}
\centering
\begin{tabular}{lrrr}
\toprule
キャンペーンタイプ & 帰属数 & 整合数 & 整合率 \\
\midrule
TypeA（大規模販促） & 18 & 10 & 0.56 \\
TypeB（ターゲット型） & 78 & 1 & 0.01 \\
TypeC（ニッチ型） & 18 & 0 & 0.00 \\
\midrule
全体 & 114 & 11 & 0.10 \\
\bottomrule
\end{tabular}
\end{table}

全体のクーポン整合率は10\%と低い．
この数字自体は本手法の限界を反映しているが，
キャンペーンタイプ別に分解すると明確な構造が浮かび上がる．
\textbf{TypeA}（大規模販促; C8, C13, C18, C26, C30）では
クーポン対象商品が3.7K--38.2K件と多く，
帰属パターンの56\%がクーポン対象商品を含む（整合率0.56）．
これは大規模販促が集団レベルのサポート変動を生み出し，
パイプラインが正しく検出していることを示唆する．
一方，\textbf{TypeB/TypeC}（ターゲット型・ニッチ型キャンペーン）は
特定世帯向けの個別施策であり，
クーポン対象商品数は200--2K件と限定的で，
集団レベルのサポート時系列には影響が表れにくい．
これらのキャンペーンに対する帰属は，
プロモーション効果ではなく季節性やトレンドなど
別の要因によるサポート変動を反映している可能性が高い．

この結果は本パイプラインの適用範囲を示唆する:
集団レベルで検出可能なサポート変動を生む\textbf{大規模イベント}に対しては
高い帰属整合性を持つが，個別世帯レベルの微細な効果の検出は
本手法の設計上の範囲外である．


% EX4 は EX1c および EX2c に統合済み
